\subsection{Reference estimators and uncertainty of the comparisons}
\label{sec:protocolo_comparacao}

On the training rows of the encoders, XGBoost and LightGBM \citep{chen2016xgboost,ke2017lightgbm} fit one regressor per coordinate by least squares to \num{58} aggregated features, such as centroids with linear and quadratic weights, quantiles, spreads, skewnesses and brightest channels. Each regressor has up to \num{3000} trees, and its training stops after \num{100} rounds without improvement on validation; the hyperparameters are those with the lowest mean \mbox{3-D} validation error among \num{40} TPE trials, \num{12} of them random \citep{akiba2019optuna}. In the run records, validation and testing used light with data augmentation (channel dropout and spurious channel activation, noise, gain variation, photon thinning), with \qty{4.1}{\percent} less light in the test set. This report re-evaluates the boosters on the clean light of the encoders; with augmentation, it reproduces the recorded values exactly. The search and the training, however, used augmented light.

The network of the previous stage \citep{vicentin2024study} is quoted rather than re-evaluated: it used a different partition (\qty{80}{\percent} and \qty{20}{\percent}), the geometry prior to the end-cap correction and, judging by its plots, the file-name frame. APEX \citep{marinho2025apex} uses a different simulation (\qty{10}{ppm} of xenon, a field cage reflectivity of \qty{79}{\percent}, light re-emitted by the photon detectors) and analytic estimators. Its resolution widths along its own axes, $\sigma_x$, $\sigma_y$ and $\sigma_z$, read from Figure~2 of the paper at \qty{0.26}{\cm} per pixel, have no published definition; the combined width $\sigma_{\mathrm{comb}}$ is their root mean square. The main comparison is at equal detected light, which does not guarantee equal position information: the mean test light of \qty{530.1}{\pe}\ corresponds to \qty{2.9}{\MeV} in APEX with the average yield, \qty{180}{\pe\per\MeV}, and to \qty{4.9}{\MeV} with the minimum, \qty{109}{\pe\per\MeV}, an energy at which even the dimmest region of APEX gives that much light, which favours APEX; the comparison is also made at equal energy, \qty{10}{\MeV}. Only $\sigma_{\mathrm{comb}}$ and $\sigma_z$ are compared: the correspondence of the transverse widths to $x$ and $y$ is not established.

The barycentre $\bar{z}$, with a quadratic calibration fitted on \num{4000} randomly drawn training events, uses the same \num{937} test events and the same photon draw. The \qty{95}{\percent} intervals come from bootstrap resampling of the test events \citep{efron1979bootstrap}, paired for the encoders and boosters, with \num{10000} shared resamples, and unpaired for the barycentre, with \num{2000}; they cover the sampling of the test set, not the variation between training runs. Because the same test events served the first comparison and the diagnosis that guided the final one (Section~\ref{sec:dificuldades}), the test errors of the encoders may be somewhat optimistic, which the intervals do not cover either. Differences, ratios and percentages are computed without rounding.
